% ============================================================
% Notes: Self-similarity, Power Laws, Fractals
% + link to anomalous diffusion / fractional calculus
% Source motivation: Herrmann, Fractional Calculus (Mandelbrot + fractals)
% ============================================================

\section{Concept of self-similarity}

\subsection{Exact self-similarity (scale invariance)}
A function/object is \textbf{self-similar} if it reproduces its form under rescaling.

\paragraph{Functional form.}
A function $f(x)$ is scale invariant if
\begin{equation}
f(\lambda x)=\lambda^{k} f(x)
\qquad \text{for all }\lambda>0,
\end{equation}
for some exponent $k\in\mathbb{R}$.

\paragraph{Solution.}
The only (nontrivial) continuous solutions are power laws:
\begin{equation}
f(x)=C\,x^{k}.
\end{equation}

\subsection{Self-similarity in stochastic processes}
A stochastic process $X(t)$ is \textbf{self-similar} with parameter $H$ (Hurst exponent) if
\begin{equation}
X(\lambda t)\ \overset{d}{=}\ \lambda^{H}X(t),
\qquad \lambda>0,
\end{equation}
where $\overset{d}{=}$ means equality in distribution.

Typical consequences:
\begin{equation}
\mathbb{E}[X(t)]\sim t^{H},
\qquad
\mathbb{E}[X^2(t)]\sim t^{2H}.
\end{equation}

\paragraph{Diffusion scaling link.}
If MSD scales as
\[
\langle X^2(t)\rangle\sim t^\alpha,
\]
then typically
\begin{equation}
\alpha = 2H.
\end{equation}

% ------------------------------------------------------------

\section{Power-law functions}

\subsection{Definition}
A \textbf{power law} is any relation of the form
\begin{equation}
y(x)=C\,x^{k},
\qquad x>0.
\end{equation}

\subsection{Key properties}
\begin{itemize}
\item \textbf{Scale invariance:}
\[
y(\lambda x)=C(\lambda x)^k=\lambda^k y(x).
\]
\item \textbf{Straight line on log--log plot:}
\begin{equation}
\log y = \log C + k\log x.
\end{equation}
So slope in a $\log y$ vs $\log x$ plot equals $k$.
\item \textbf{No intrinsic scale:} power laws do not define a characteristic length/time.
\end{itemize}

\subsection{Power-law tails (heavy tails)}
A probability density can have a power-law tail:
\begin{equation}
p(x)\sim \frac{1}{|x|^{1+\mu}}
\qquad (|x|\to\infty).
\end{equation}

Moments may diverge:
\begin{equation}
\mathbb{E}[|X|^q] < \infty \iff q<\mu.
\end{equation}

Example: if $0<\mu<2$, then variance diverges.

\subsection{Why power laws appear}
Power laws typically arise from:
\begin{itemize}
\item scale invariance / self-similarity
\item multiplicative processes
\item critical phenomena (no characteristic scale)
\item fractal geometry
\item long-memory kernels (e.g. $t^{-\alpha}$)
\end{itemize}

% ------------------------------------------------------------

\section{Fractals}

\subsection{Basic idea}
A \textbf{fractal} is a set/geometry with:
\begin{itemize}
\item structure at many scales (self-similarity),
\item non-integer (fractional) dimension,
\item often power-law scaling of measured quantities.
\end{itemize}

\subsection{Self-similar fractals and fractal dimension}
Suppose a set is made of $N$ self-similar copies of itself,
each scaled by a factor $r$ (with $0<r<1$).

Then the \textbf{similarity dimension} $D$ satisfies
\begin{equation}
N = r^{-D}
\qquad\Longleftrightarrow\qquad
D=\frac{\log N}{\log(1/r)}.
\end{equation}

\paragraph{Examples.}
\begin{itemize}
\item \textbf{Cantor set:} $N=2$, $r=1/3$
\[
D=\frac{\log 2}{\log 3}\approx 0.6309.
\]
\item \textbf{Koch curve:} $N=4$, $r=1/3$
\[
D=\frac{\log 4}{\log 3}\approx 1.2619.
\]
\item \textbf{Sierpiński triangle:} $N=3$, $r=1/2$
\[
D=\frac{\log 3}{\log 2}\approx 1.585.
\]
\end{itemize}

\subsection{Box-counting (Minkowski) dimension}
Cover the set $S$ with boxes of side length $\varepsilon$.
Let $N(\varepsilon)$ be the minimal number needed.

The \textbf{box-counting dimension} is
\begin{equation}
D_B = \lim_{\varepsilon\to 0}\frac{\log N(\varepsilon)}{\log(1/\varepsilon)},
\end{equation}
if the limit exists.

\subsection{Scaling laws on fractals}
On fractals, measured quantities often scale with non-integer exponents.

Example: ``mass'' or number of points inside radius $R$:
\begin{equation}
M(R)\sim R^{D},
\end{equation}
where $D$ is the fractal dimension.

\subsection{Random walks on fractals (important physics link)}
Diffusion on fractal/disordered structures often becomes anomalous:
\begin{equation}
\langle r^2(t)\rangle \sim t^\alpha,
\qquad \alpha\neq 1.
\end{equation}

This is one route to \textbf{fractional} models of transport.

% ------------------------------------------------------------

\section{Connection to fractional calculus (Herrmann motivation)}

Herrmann notes that Mandelbrot's fractal geometry (1980s) strongly attracted physicists
and produced major interest in fractional Brownian motion and anomalous diffusion,
which are key motivations for fractional calculus methods in physics. :contentReference[oaicite:1]{index=1}

\subsection{Why fractals naturally lead to fractional operators}
Heuristic links:
\begin{itemize}
\item \textbf{Fractal geometry} $\Rightarrow$ effective non-integer dimension / scaling.
\item \textbf{Power-law memory kernels} $\Rightarrow$ nonlocal time operators.
\item \textbf{Long-range jump statistics} $\Rightarrow$ fractional Laplacians in space.
\end{itemize}

Thus, fractional derivatives are natural mathematical tools to encode
\textbf{scale-free} and \textbf{nonlocal} behavior.

% ------------------------------------------------------------

\section{High-yield summary}
\begin{itemize}
\item \textbf{Self-similarity} = invariance under scaling.
\item \textbf{Scale invariance equation:} $f(\lambda x)=\lambda^k f(x)$.
\item \textbf{Power laws} are the unique continuous solutions: $f(x)=Cx^k$.
\item \textbf{Power laws} appear as straight lines on log--log plots.
\item \textbf{Fractals} = self-similar sets with non-integer dimension.
\item \textbf{Similarity dimension:} $D=\log N/\log(1/r)$.
\item \textbf{Box-counting dimension:} $D_B=\lim_{\varepsilon\to 0}\log N(\varepsilon)/\log(1/\varepsilon)$.
\item \textbf{Physics link:} fractals $\Rightarrow$ anomalous diffusion $\Rightarrow$ fractional calculus models. :contentReference[oaicite:2]{index=2}
\end{itemize}
